\section{The Three Pillars of Learning}
\begin{frame}{The Three Pillars of Learning}
    \centering
    % Scaled down to 85% to fit better
    \begin{tikzpicture}[
        scale=0.85, transform shape, 
        node distance=1.0cm and 0.3cm,
        box/.style={
            rectangle, rounded corners, draw=none, 
            text width=4.0cm, minimum height=1.8cm, 
            align=center, font=\small
        },
        title/.style={font=\bfseries\large},
        % Arrow style removed
    ]

        % --- Column 1: The Model ---
        \node[title, blue] (h_title) at (0,0) {1. The Model ($f_\theta$)};
        \node[box, fill=blue!10, below=0.2cm of h_title] (h_box) {
            \textbf{Design Choice:}\\
            ML (Linear, Kernel, Trees) \\ \& DL (MLP, CNN)
            \vspace{0.5em}
            
            \textbf{Depends on:}\\
            Data Type (Tabular, image) \\ \& Complexity
        };

        % --- Column 2: Loss Function ---
        \node[title, red] (l_title) at (4.5,0) {2. Loss Function ($J$)};
        \node[box, fill=red!10, below=0.2cm of l_title] (l_box) {
            \textbf{Design Choice:}\\
            MSE vs. Cross-Entropy
            \vspace{0.5em}
            
            \textbf{Depends on:}\\
            The Task\\
            (Regression vs. Classification)
        };

        % --- Column 3: Optimizer (Dimmed) ---
        \node[title, green!60!black] (o_title) at (9.0,0) {3. Optimizer};
        \node[box, fill=green!10, below=0.2cm of o_title] (o_box) {
            \textbf{Design Choice:}\\
            Gradient Descent, Adam
            \vspace{0.5em}
            
            \textbf{Depends on:}\\
            Stability \& Convergence Speed\\
            \footnotesize{(Covered later in Deep Learning)}
        };

        % Arrows removed as requested

    \end{tikzpicture}

    \vspace{1em}

    \begin{block}{}
        Today we focus on \textbf{1 \& 2}. \\
        % Let's start by discussing the loss functions first.
    \end{block}
\end{frame}


\begin{frame}{The Three Pillars of Learning}
    \centering
    % Scaled down to 85% to fit better
    \begin{tikzpicture}[
        scale=0.85, transform shape, 
        node distance=1.0cm and 0.3cm,
        box/.style={
            rectangle, rounded corners, draw=none, 
            text width=4.0cm, minimum height=1.8cm, 
            align=center, font=\small
        },
        title/.style={font=\bfseries\large},
        % Arrow style removed
    ]

        % --- Column 1: The Model (Dimmed) ---
        \node[title, blue, opacity=0.4] (h_title) at (0,0) {1. The Model ($f_\theta$)};
        \node[box, fill=blue!10, text=gray, opacity=0.5, below=0.2cm of h_title] (h_box) {
            \textbf{Design Choice:}\\
            ML (Linear, Kernel, Trees) \\ \& DL (MLP, CNN)
            \vspace{0.5em}
            
            \textbf{Depends on:}\\
            Data Type (Tabular, image) \\ \& Complexity
        };

        % --- Column 2: Loss Function (Active) ---
        \node[title, red] (l_title) at (4.5,0) {2. Loss Function ($J$)};
        \node[box, fill=red!10, below=0.2cm of l_title] (l_box) {
            \textbf{Design Choice:}\\
            MSE vs. Cross-Entropy
            \vspace{0.5em}
            
            \textbf{Depends on:}\\
            The Task\\
            (Regression vs. Classification)
        };

        % --- Column 3: Optimizer (Dimmed) ---
        \node[title, green!60!black, opacity=0.4] (o_title) at (9.0,0) {3. Optimizer};
        \node[box, fill=green!10, text=gray, opacity=0.5, below=0.2cm of o_title] (o_box) {
            \textbf{Design Choice:}\\
            Gradient Descent, Adam
            \vspace{0.5em}
            
            \textbf{Depends on:}\\
            Stability \& Convergence Speed\\
            \footnotesize{(Covered later in Deep Learning)}
        };

        % Arrows removed as requested

    \end{tikzpicture}

    \vspace{1em}

    \begin{block}{}
        Today we focus on \textbf{1 \& 2}. \\
        Let's start by discussing the loss functions first.
    \end{block}
\end{frame}


% =========================
% 2. Supervised Learning
% =========================

% \section{Supervised Learning Algorithms}

% \subsection{Problem Formulation}

% \begin{frame}{Supervised Learning: Problem Setup}
%     In supervised learning, we are provided with a labeled dataset $\{(x_i, y_i)\}_{i=1}^{n}$, where each input $x_i \in \mathbb{R}^d$ is paired with a target $y_i$. The objective is to learn a function $f: \mathbb{R}^d \to \mathbb{R}$ (or discrete space) that generalizes well to unseen data, such that $f(x_i) \approx y_i$ for all $i$.
    
%     \vspace{1em}
    
%     Supervised tasks divide into two primary categories based on the nature of $y_i$:
    
%     \begin{columns}[T]
%         \begin{column}{0.48\textwidth}
%             \textbf{Regression Tasks}
%             \begin{itemize}
%                 \item $y_i$ is continuous ($y_i \in \mathbb{R}$)
%                 \item Examples: Predicting house prices or temperatures
%                 \item Common loss: Mean Squared Error (MSE), which measures average squared deviation
%             \end{itemize}
%         \end{column}
        
%         \begin{column}{0.48\textwidth}
%             \textbf{Classification Tasks}
%             \begin{itemize}
%                 \item $y_i$ is categorical ($y_i \in \{1, \dots, K\}$)
%                 \item Examples: Identifying spam emails or medical diagnoses
%                 \item Common loss: Cross-Entropy, which quantifies probabilistic mismatch
%             \end{itemize}
%         \end{column}
%     \end{columns}
% \end{frame}

% =================================================================
% SECTION 1: SUPERVISED LEARNING FRAMEWORK
% =================================================================

\section{Supervised Learning}



\begin{frame}{Supervised Learning: The Setup}
    \begin{block}{The Objective}
        Given a labeled dataset $\mathcal{D} = \{(x_i, y_i)\}_{i=1}^{n}$, we aim to learn a mapping function $f$ that generalizes to unseen data:
        \begin{equation*}
            \hat{y} = f(x; \theta) \approx y
        \end{equation*}
        \small where $\theta$ represents the learnable model parameters.
    \end{block}

    \vspace{1.5em}

    \centering
    \begin{tikzpicture}[node distance=2cm, auto, thick, >=stealth]
        % Nodes
        \node[draw, circle, fill=gray!10, minimum size=1cm] (input) {$x_i$};
        \node[draw, rectangle, rounded corners, fill=blue!10, minimum width=2.5cm, minimum height=1.2cm, right=of input] (model) {Model $f_\theta$};
        \node[draw, circle, fill=green!10, minimum size=1cm, right=of model] (output) {$\hat{y}_i$};
        \node[draw, circle, fill=gray!10, minimum size=1cm, below=0.8cm of output] (target) {$y_i$};
        
        % Edges
        \draw[->] (input) -- node[above, font=\scriptsize] {Features} (model);
        \draw[->] (model) -- node[above, font=\scriptsize] {Prediction} (output);
        \draw[<->, dashed, red] (output) -- node[right, font=\scriptsize] {Error / Loss} (target);
        
        % Annotation
        \node[below=0.5cm of model, text width=5cm, align=center, font=\itshape\small] 
            {Learning involves adjusting $\theta$ to minimize this error across all $i$.};
    \end{tikzpicture}
\end{frame}



% % --- SLIDE 1: THE MATHEMATICAL SETUP ---
% \begin{frame}{1.1 Supervised Learning: The Setup}
%     \begin{block}{The Objective}
%         Given a training set $\mathcal{D} = \{(\bx^{(i)}, y^{(i)})\}_{i=1}^N$, learn a mapping function $f_\theta$ such that:
%         $$ f_\theta(\bx^{(i)}) \approx y^{(i)} $$
%         for unseen data (Generalization).
%     \end{block}

%     \vspace{1.5em}
%     \centering
%     % Visualizing the Function Mapping
%     \begin{tikzpicture}[node distance=2cm, auto, thick]
%         % Nodes
%         \node[draw, circle, fill=gray!10, minimum size=1cm] (input) {$\bx$};
%         \node[draw, rectangle, fill=blue!10, minimum width=2.5cm, minimum height=1.2cm, right=of input] (model) {Model $f_\theta$};
%         \node[draw, circle, fill=green!10, minimum size=1cm, right=of model] (output) {$\hat{y}$};
%         \node[draw, circle, fill=gray!10, minimum size=1cm, below=0.8cm of output] (target) {$y$};
        
%         % Edges
%         \draw[->] (input) -- node[above] {Features} (model);
%         \draw[->] (model) -- node[above] {Prediction} (output);
%         \draw[<->, dashed, red] (output) -- node[right] {Error (Loss)} (target);
%     \end{tikzpicture}
% \end{frame}

\subsection{Regression}

% --- SLIDE 2: REGRESSION TASKS ---
